\chapter{API für Struktur-Topologieoptimierung}\label{api:topology}
Version 1.6.0 ergänzt die Module \path{topology},
\path{topology_optimization}, \path{topology_study} und \path{topology_geometry}.
Alle Dichte-/Gradientenvektoren haben Form \((n_{\rm Zellen},)\), die
Eingabezellfolge bleibt erhalten. \path{LoadCase} enthält konstante SI-Lasten
und homogene Dirichlet-DOFs; \path{TopologyProblem} validiert Tet4, positive
Tensoren, Gewichte, Radius sowie feste Voll-/Leergebiete.
\path{TopologyProblem.from_models} verwendet den additiven Snapshot
\path{Model.elasticity_inputs()}; Material- und Lastparser bleiben gemeinsam.
\path{evaluate} liefert unveränderliche Dichten, Compliance in J,
Volumenfraktion und vollständige Ableitungen sowie je ein \path{ElasticResult}.

\path{OptimizationOptions} begrenzt Iterationen/Auswertungen und definiert die
skalierte SLSQP-Toleranz. \path{optimize_topology} bewahrt den besten zulässigen
Zustand bei Budget-/Solverstopp. Status sind \path{converged},
\path{iteration_limit}, \path{evaluation_limit} und \path{solver_failed}.
Ein ungültiger Anfangszustand wirft \path{ValidationError}.
\begin{lstlisting}
study = TopologyStudy.from_file("study.toml")
result = study.run("new-output")
record = read_topology_result("new-output")
export_reanalyses(study.problem, result.final,
                  "new-output/geometry")
\end{lstlisting}
TOML und Ergebnis verwenden eigene Schema-Version 1. Eingaben enthalten
vollständige effektive SI-Daten, Solverparameter und Quellhashes;
\path{RunProvenance} bindet sie deterministisch. JSON/VTU erhalten eine
SHA-256-Inventardatei. Vorhandene Ausgabeordner werden abgelehnt.
\path{reanalyze_geometry(..., connect=False)} prüft die reine Schwelle;
\path{connect=True} ergänzt ausdrücklich protokollierte Verbindungszellen.
\path{GeometryResult} enthält das volle diskrete Netz, Originalzellzeilen,
Zustände, Nachgiebigkeit, Volumen und rohe von-Mises-Spannung. Kein CAD,
keine Spannungs-/Fertigungsfreigabe und keine implizite Änderung alter Formate.
